\label{sec:legal}

\subsection{Administrative Subdivision Integrity in Constitutional Doctrine}

The constitutional authority for treating administrative subdivisions as communities
of interest begins with \textit{Reynolds v.\ Sims}, 377 U.S.\ 533 (1964)~\citep{reynolds1964}.
While the primary holding of \textit{Reynolds} was that legislative seats must be
apportioned on a population basis, Chief Justice Warren's majority opinion expressly
preserved the legitimacy of administrative subdivision preservation as a secondary
criterion:
\begin{quote}
  ``A State may legitimately desire to maintain the integrity of various political
  subdivisions, insofar as possible, and provide for compact districts of contiguous
  territory in apportioning its legislature.
  So long as the divergences from a strict population standard are based on legitimate
  considerations incident to the effectuation of a rational state policy, some
  deviations from the equal-population principle are constitutionally permissible.''
\end{quote}
The Court's acknowledgment that ``integrity of various political subdivisions'' is a
legitimate consideration is the constitutional anchor for zone co-membership weights.
Counties, municipalities, townships, and school districts are all political subdivisions
in the constitutional sense; the zone co-membership score directly measures how many
such subdivisions two tracts share.

\subsection{School District Integrity and Milliken v.\ Bradley}

\textit{Milliken v.\ Bradley}, 418 U.S.\ 717 (1974)~\citep{milliken1974},
arose from a desegregation case in the Detroit metropolitan area.
The Supreme Court declined to order a cross-district busing remedy absent evidence
that the school district lines had been drawn for segregatory purposes.
Justice Stewart's concurrence, joined by the majority, treated school district lines
as presumptively legitimate community boundaries:
\begin{quote}
  ``No single tradition in public education is more deeply rooted than local control
  over the operation of schools; local autonomy has long been thought essential both
  to the maintenance of community concern and support for public schools and to
  quality of the educational process.''
\end{quote}
While \textit{Milliken} addressed desegregation remedies rather than legislative
redistricting directly, its reasoning applies with equal force: school districts
represent genuine local communities organized around the shared civic enterprise of
public education, and disrupting school district boundaries imposes real costs on
those communities.

In the redistricting context, this reasoning supports the proposition that a district
plan that splits a school district has imposed a cost on a genuine community of
interest.
The zone co-membership formula captures this cost by making the cut between
co-school-district tracts more expensive than the cut between tracts in different
school districts.

\subsection{The Property Tax Fiscal Bond}

The property tax is the financial foundation of local government in the United States.
It funds public schools (36\% of school revenue nationally~\citep{nces2022}),
fire protection, local roads, libraries, and parks.
More importantly for redistricting purposes, it creates a direct fiscal bond between
neighbors: two households in the same school district are co-investors in the same
educational enterprise.
When the school district calls a bond election to finance a new gymnasium, every
property owner in the district pays if the measure passes.
This creates a community of shared fiscal interest that is as tangible and verifiable
as any community of interest recognized in redistricting law.

The property tax fiscal bond operates through county subdivisions in New England states.
A Massachusetts town, a Vermont town, or a Connecticut town is the unit that levies
the local property tax, maintains local roads, operates the town hall, and holds
the annual town meeting at which residents vote directly on the budget.
Two census tracts in the same Massachusetts town are co-owners of a civic enterprise
in a direct and literal sense; tracts split between different towns are not.

Zone co-membership scores naturally capture this fiscal bond.
In Massachusetts, the county subdivision (town) dimension of the zone score is the
dominant signal: towns are small, well-defined, and politically salient.
In North Carolina, the school district dimension is more important: school districts
are large but cross county boundaries, making them a more meaningful unit than the
county.
The formula weights all available zone types equally, which is appropriate because
the relative importance of zone types varies by state.

\subsection{Electric Utility Rate-Payer Community}

State public utility commission law provides an independent legal foundation for
electric utility service territory as a community of interest.
When a utility files a rate case before the public utility commission, all customers
in the service territory have standing as a class.
The commission may consolidate customer objections, appoint a consumer advocate,
and conduct hearings in the service territory.
The rate-payers of a service territory are, in state regulatory law, a recognized
interest group with collective legal standing.

This community relationship is reinforced by the exclusive franchise nature of
electric utility service.
Unlike school districts, which sometimes overlap or where families can exercise
choice, electric utility service territories are exclusive: a tract is served by
exactly one retail electric utility and cannot choose a different provider.
This exclusivity makes the rate-payer community bond non-voluntary and durable in
a way that other community memberships are not.

The EIA Form 861 data capturing this bond is collected annually, verified by
federal statute (Energy Policy Act of 2005), and publicly available at no cost.
It is the most reliable and uniformly available zone type in the M.6 formula.

\subsection{The Shaw v.\ Reno Neutrality Requirement}

\textit{Shaw v.\ Reno}, 509 U.S.\ 630 (1993)~\citep{shaw1993}, requires that
race not be the predominant factor in drawing district lines.
The zone co-membership formula satisfies this requirement because:
\begin{enumerate}
  \item No racial composition data enters the zone assignment.
    Zone membership is determined entirely by the spatial boundaries of school
    districts, county subdivisions, utility service territories, and other
    administrative zones.
    These boundaries are determined by historical administrative decisions,
    not by racial composition.
  \item The zone co-membership score does not weight zone types differently based
    on the racial composition of tract residents.
    School district co-membership is weighted the same way regardless of whether
    the tracts are predominantly white or predominantly minority.
  \item The formula is the same for all tract pairs: $\mathrm{zone\_score}(u,v) =
    |\text{shared zones}| / |Z|$.
    There is no race-specific adjustment.
\end{enumerate}

The fact that administrative zones may correlate with race---as virtually any
geographic feature does---does not make the zone co-membership signal racially
motivated under the Shaw-Miller standard.
The Supreme Court in \textit{Miller v.\ Johnson}, 515 U.S.\ 900 (1995)~\citep{miller1995},
made clear that the relevant question is whether race was the \emph{predominant factor},
not whether racial effects existed.
A zone co-membership weight that happens to preserve majority-minority school
districts because those districts are genuine administrative communities satisfies
the constitutional standard.

\subsection{State Commission Acceptance}

Administrative subdivision integrity has been accepted as a communities of interest
criterion by redistricting commissions in California, Colorado, and Arizona.
California's Citizens Redistricting Commission has cited school district preservation
as a specific operationalization of the communities of interest criterion in its
2011 and 2021 proceedings.
Colorado's independent commission has used county and municipal integrity as a proxy
for communities of interest.
Arizona's IRC has accepted expert testimony about administrative zone co-membership
as evidence that communities of interest have been preserved.

The zone co-membership formula provides a quantitative, replicable version of the
qualitative judgments that these commissions have made manually.
Rather than relying on a commissioner's judgment that ``this neighborhood is clearly
a community of interest,'' the formula computes a score that any party can verify
using the same public data.
This quantitative defensibility is the primary advantage of the zone co-membership
approach over qualitative community identification.
